\section{Strong-Weak Representative Involution}
\label{sec:strong_weak_involution}

\paragraph{Physical idea.}
In many closed mean-force parameterizations, one scalar controls the crossover between two asymptotic
languages: a weak-influence expansion in powers of \(\chi\), and a strong-influence expansion in
powers of \(1/\chi\).
The key duality is therefore a reflection of the crossover coordinate,
\begin{equation}
  y=\log\chi \longmapsto -y,
\end{equation}
which is equivalent to
\begin{equation}
  \chi\longmapsto \chi^{-1}.
\end{equation}

\paragraph{Representative map at fixed \(\beta\).}
For one-channel scaling,
\begin{equation}
  \chi(\beta,g,\vartheta)=g^2\chi_0(\beta,\vartheta),
\end{equation}
define
\begin{equation}
  g_\star(\beta,\vartheta):=\chi_0(\beta,\vartheta)^{-1/2},
  \qquad
  g^\vee:=\frac{g_\star^2}{g}=\frac{1}{\chi_0 g}.
\end{equation}
Then
\begin{equation}
  \chi(\beta,g^\vee,\vartheta)=\chi(\beta,g,\vartheta)^{-1},
  \qquad
  (g^\vee)^\vee=g.
\end{equation}
So each point has a partner across the self-dual contour \(\chi=1\).

\paragraph{Control transfer, not state equality.}
The dual map does \emph{not} claim
\(Q(\beta,g,\vartheta)=Q(\beta,g^\vee,\vartheta)\) pointwise.
Instead, it states that the hard strong-branch expansion at \((\beta,g)\) can be reconstructed from
the weak-branch expansion of the same functional family at \((\beta,g^\vee)\).
This is a computational statement about asymptotic control.

\paragraph{Concrete example.}
For
\begin{equation}
  \gamma(\chi)=\frac{\tanh\chi}{\chi},
\end{equation}
one has
\begin{align}
  \gamma(\chi) &\sim 1-\frac{\chi^2}{3}+O(\chi^4), && \chi\to0, \\
  \gamma(\chi) &\sim \frac{1}{\chi}, && \chi\to\infty.
\end{align}
Under \(\chi\mapsto\chi^{-1}\), the weak expansion variable at the dual representative is exactly
the strong expansion variable of the original representative.

\paragraph{Geometric reading.}
In \((\beta,g)\)-space, \(\chi=1\) is the self-dual manifold.
The map \(g\leftrightarrow g^\vee\) at fixed \(\beta\) is a reflection across this manifold in the
coordinate \(y=\log\chi\).
The normal-duality construction later in the paper generalizes this to simultaneous
\((\beta,g)\)-motion along local normals.

\paragraph{Formal statement.}
A theorem-level formulation (including admissibility conditions and proof) is given in the Appendix.
